
\section{Modeling details}

\subsection{Tuning the hyperparameters of the training objective}
\label{sec:training_objective_hyperparams}
We find it practical to scale the reconstruction loss by a dataset-specific scalar $\alpha$ since
the MSE loss varies in average magnitude across datasets. When training H-VAE, we find it
beneficial to keep $\beta = 0$ for a few epochs ($E_\text{rec}$) so that the model can learn to
perform meaningful reconstructions, and then linearly increasing it to the desired value for
$E_\text{warmup}$ epochs to add structure to the latent space, an approach first used by
\citet{bowman_generating_2016}.

\subsection{Data normalization}
\label{sec:data_normalization}
As per standard machine learning practice~\cite{shanker_effect_1996}, we normalize the data. We do
this by dividing each tensor by the average square-root total norm of the training tensors,
analogously to the Signal Norm. This strategy puts the target values on a similar scale as the
normalized activations learned by the network, which we speculate to favor gradient flow.

\subsection{Implementation details}
Without loss of generality, we use real spherical harmonics for implementation of H-(V)AE. 
We leverage \verb|e3nn|~\cite{geiger_e3nn_2022}, using their computation of the real spherical
harmonics and their Clebsch-Gordan coefficients.\\
\\
In our code, we offer the option to use the Full Tensor Product instead of the ETP. Specifically,
at each block we allow the users to specify whether to compute the Tensor Product element-wise or
fully-connected within the multiplicity, and whether to compute using efficient or fully connected
degree mixing.

\subsection{Pairwise invariant reconstruction loss}
\label{sec:pairwise_inv_rec_loss}
To reconstruct a signal within an equivariant model it is desirable to have a \textit{pairwise
invariant} reconstruction loss, i.e., a loss $\mathcal{L}_{rec}$ such that $\mathcal{L}_{rec}(\vx,
\vy) = \mathcal{L}_{rec}(\mD(\mathfrak{g})\vx, \mD(\mathfrak{g})\vy)$ where $\mD$ is the
representation of the group element $\mathfrak{g}$ acting on the space that $x$ and $y$ inhabit
(e.g. a rotation matrix if $\vx$ and $\vy$ are vectors in Euclidean 3D space, or a degree-$\ell$
wigner-D matrix if $\vx$ and $\vy$ are degree-$\ell$ vectors). This  property is necessary for the
model to remain equivariant, i.e., given that the network is agnostic to the transformation of the
input under group operation $\vx \to \mD(\mathfrak{g})\vx$ by producing a similarly transformed
output $\vy \to \mD(\mathfrak{g})\vy$, we want the reconstruction loss to be agnostic to the same
kind of transformation as well\\
\\
The MSE loss is pairwise invariant for any degree-$\ell$ feature on which SO(3) acts via the
$\ell$'s Wigner-D matrix. Consider two degree-$\ell$ features $\vx_{\ell}$ and $\vy_{\ell}$ acted
upon by a Wigner-D matrix $D_{\ell}(\mathfrak{g})$ parameterized by rotation $\mathfrak{g}$ (we
drop the $\mathfrak{g}$ and $\ell$ indexing for clarity):
\begin{align}
 \nonumber    \text{MSE}(\mD\vx, \mD\vy) &= (\mD\vx - \mD\vy)^{T} (\mD\vx - \mD\vy) \\
  \nonumber                       &= (\mD(\vx - \vy))^{T} (\mD(\vx - \vy)) \\
  \nonumber                       &= (\vx - \vy)^{T}\mD^{T} \mD(\vx - \vy) \\
  \nonumber                       &= (\vx - \vy)^{T} (\vx - \vy) \qquad \qquad \text{since Wigner-D matrices are unitary} \\ 
  \nonumber                       &= \text{MSE}(\vx, \vy)\\
\end{align}
Since the MSE loss is pairwise invariant for every pair of degree-$\ell$ features, it is thus
pairwise invariant for pairs of steerable tensors composed via direct products of steerable
features.

\subsection{Cosine loss}
\label{sec:cosine_loss_appendix}
We use the metric {\em Cosine loss} to measure a model's reconstruction ability. Cosine loss is a
normalized dot product generalized to operate on pairs of steerable tensors (akin to cosine
similarity), and modified to be interpreted as a loss (see Sec.~\ref{sec:cosine_loss_appendix} for
details),
\begin{equation}
\begin{split}
    \text{Cosine}(\vx, \vy) = 1 - \frac{\vx \odot \vy}{\sqrt{(\vx \odot \vx)(\vy \odot \vy)}}, \\
    \text{with}\quad \vx \odot \vy = \sum_{\ell'} (\vx_{\ell'} \otimes_{cg} \vy_{\ell'})_{\ell = 0}
\end{split}
    \label{eqn:cosine_loss}
\end{equation}
Importantly, unlike MSE, which depends on the characteristics of the data (e.g., the size of the
data tensors),  Cosine loss is dimensionless and therefore, interpretable and comparable across
datasets.  A measure with these characteristics is practically useful for training of a network
because it provides an estimate for how much better the reconstructions can get if the network’s
hyperparameters were to be further optimized. For example, looking at the Cosine loss in
Table~\ref{table:mse_and_cosine}, we see that our model trained on Shrec17 (best Cosine = 0.130) is
not as well optimized as our model trained on MNIST (best Cosine = 0.017). Using MSE, the trend is
reversed (1.8$\times 10^{-3}$ vs. 6.7$\times 10^{-3}$), since the scale of MSE depends on the size
of the irreps of the data (Fig.~\ref{fig:cosine_vs_mse}). Nonetheless, as a measure of a model's
reconstruction ability, Cosine loss correlates almost perfectly with MSE for a given dataset and
network, especially in the mid-to-low reconstruction quality regime (SpearmanR~$= 0.99$,
Fig.~\ref{fig:cosine_vs_mse} and Table~\ref{table:mse_and_cosine}).It should also be noted that 
since  Cosine loss ignores the relative norms of data features, it is unable to reconstruct norms
and thus, not suitable as a training objective. \\


\noindent \textbf{Proof: Invariance of  Cosine loss.}
The generalized dot product $\odot$ from Eq.~\ref{eqn:cosine_loss} is pairwise invariant in the
same way that the dot product between two 3D vectors depends only on their relative orientations
but not the global orientation of the whole two-vector system. Therefore, the whole Cosine loss
expression is pairwise invariant, since all of its components are pairwise invariant.\\


\section{Using a non-equivariant decoder}
\label{sec:mlp_decoder_appendix}
{\em Winter et al}~\cite{winter_unsupervised_2022} propose to construct group-equivariant
autoencoders by using an equivariant encoder that learns an invariant embedding and a group
element, and an unconstrained decoder which uses the invariants alone to reconstruct each datapoint
in the "canonical" form, before applying the learned group action in the output space. By contrast,
for SO(3) we propose to use an equivariant decoder, whereby the learned group element is fed as
input to the decoder. Such ``unconstrained decoder" procedure can in principle be merged with our
equivariant encoder and Fourier-space approach in two ways. For each, we argue in favor of using
our equivariant decoder.\\
\\
\textbf{1) Reconstructing the Fourier coefficients of the data.} To apply the learned group element
on the decoder’s output, the Wigner-D matrices for the data’s irreps need to be computed from the
group element. Then, the Wigner-D matrices can be used to “rotate” the tensor. This has to be done
on-the-fly, and it can be done quickly using functions provided in the e3nn
package~\cite{geiger_e3nn_2022} and by smartly vectorizing operations. We implemented this
procedure by using a simple Multi-Layer Perceptron with SiLU non-linearities as a decoder. By using
e3nn to compute Wigner-D matrices in batches, and by clever construction of tensor multiplications
such that runtime scales linearly with $\ell_{max}$ and is constant with regards to multiplicity
and batch size, we achieve models that run with comparable speed to those using our equivariant
decoder, and have comparable performance on MNIST (Table~\ref{table:mlp_decoder_comparison}). Given
the empirical similarities we observe, though on a limited use case, we favor the simplicity and
elegance of our equivariant decoder. “Simplicity” because we construct the decoder to be symmetric
to the encoder, thus endowing it automatically with similar representational power and without the
need to tune an architecture made with different base components. Furthermore, we highlight that
our method generates intermediate equivariant representations in the decoder, rather than
intermediate invariant representations. These intermediate equivariant representations may be of
interest to study in and of themselves.\\
\\
\textbf{2) Reconstructing the data in real space.} In this case, we do not have to compute Wigner-D
matrices on-the-fly, since the learned frame can be used directly in the output space as a rotation
matrix. However, since the encoder only sees a truncated Fourier representation of the data, which
is by construction lossy, while the loss is computed over fine-grained real-space, this model might
be too difficult to train. We suspect this would make the model akin to a denoising
autoencoder~\cite{vincent_extracting_2008} and it might be interesting to analyze, but that would
be beyond the scope of this paper. To avoid the denoising effect, we could learn to reconstruct
data in real space after an Inverse Fourier Transform (IFT). However, computing the IFT on-the-fly
is very expensive and requires a discretization of the input space, to the point of being
prohibitive for point clouds. This is not a bottleneck for the Forward Fourier Transform if the
cloud is parameterized by Dirac-Delta distributions, i.e., for point clouds, as the integral can be
computed exactly (Eq.~\ref{eq:ZFT}).


\begin{table}[h]
\begin{center}
\resizebox{0.9\textwidth}{!}{
\begin{tabular}{l c | c | c c | c c | c c}
    \bf Method & \bf z & \bf Speed & \bf MSE & \bf Cosine & \bf Purity & \bf V-meas. & \bf
    \newlinecell{LC Class.\\Acc.} & \bf \newlinecell{KNN Class.\\Acc.} \\
    \hline
    \hline
    H-AE NR/R                       & 16 & \textbf{1.0} & $1.2 \times 10^{-3}$ & 0.031 &
    \textbf{0.66} & \textbf{0.55} & 0.850 & 0.902 \\
    H-AE unconst. decoder NR/R     & 16 & 1.3 & $\mathbf{1.1 \times 10^{-3}}$ & \textbf{0.028} &
    0.62 & 0.53 & 0.838 & \textbf{0.907} \\
    \hline
    H-VAE NR/R                      & 16 & \textbf{1.0} & $\mathbf{2.8 \times 10^{-3}}$ &
    \textbf{0.068} & \textbf{0.73} & \textbf{0.60} & \textbf{0.878} & \textbf{0.897} \\
    H-VAE unconst. decoder NR/R    & 16 & 1.3 & $2.9 \times 10^{-3}$ & 0.073 & 0.68 & 0.56 & 0.850
    & 0.872 \\
\end{tabular}
}
    \caption{\textbf{Performance comparison between our H-(V)AE and a H-(V)AE with
    ref.~\cite{winter_unsupervised_2022}'s non-equivariant decoder formulation, on the
    MNIST-on-the-sphere dataset.}}
    \medskip
    \small
    \justifying
    The non-equivariant decoders are constructed as simple MLPs with SiLU non-linearities, with the
    following hidden layer sizes: [32,64,128,160,256]. We keep the number of parameters
    approximately the same to make model comparison fair, but we do not tune the architecture of
    the invariant decoders. All other training details are kept the same
    (Sec.~\ref{sec:exp_details}).
\label{table:mlp_decoder_comparison}
\end{center}
\end{table}


\section{Experimental details}
\label{sec:exp_details}

\subsection{Architecture specification}
We describe model architectures as follows. We specify the number of blocks $B$, which is the same
for the encoder and the decoder. We specify two lists, (i) DegreesList which contains the maximum
degree $\ell_{{max},b}$ of the output of each block $b$, and (ii) MultList, containing the
multiplicity sizes $C_{b}$, of each block $b$. These lists are in the order as they appear in the
encoder, and are reversed for the decoder. When it applies, we specify the output multiplicity of
the initial linear projection $C_\text{init}$. As noted in the main text, we use a fixed formula to
determine $\ell_{{max},b}$, but we specify it for clarity.


\subsection{MNIST on the sphere}
\label{sec:mnist_details}
\textbf{Model architectures.}
For models with invariant latent space size $z = 16$, we use 6 blocks, $\text{DegreesList} = [10,
10, 8, 4, 2, 1]$ and $\text{MultList} = [16, 16, 16, 16, 16, 16]$, with a total of 227k
parameters.\\
For models with invariant latent space size $z = 120$, we use 6 blocks, $\text{DegreesList} = [10,
10, 8, 4, 2, 1]$ and $\text{MultList} = [16, 16, 16, 32, 64, 120]$, with a total of 453k
parameters.\\
\\
\textbf{Training details.}
We keep the learning schedule as similar as possible for all models. We use $\alpha = 50$. We train
VAE models for $80$ epochs using the Adam optimizer~\cite{kingma_adam_2017} with
default parameters, a batch size of $100$, and an initial learning rate of $0.001$, which we
decrease exponentially by one order of magnitude over $25$ epochs. We tune $\beta$
according to the procedure in outlined below, $E_\text{rec} = 25$ and $E_\text{warmup} = 35$.
We instead train AE models for 50 epochs using the "reduce learning rate on plateau"
scheme, reducing the learning rate by a factor of 5 if the validation loss has not improved for
over one epoch. We utilize the model with the lowest loss on validation data, only after the end
of the warmup epochs for VAE models. Training took ~$\sim 3.3$ minutes per epoch on a
single NVIDIA A40 GPU for each model.\\
\\

\textbf{Tuning the regularization strength $\beta$ with the help of conditioning.}
As outlined in Section~\ref{sec:mnist_results}, we select optimal values of $\beta$ by optimizing
the expected sample quality of the generated samples. We define samples to be of high quality if
they 1) can be correctly classified by a classifier trained on real data, and 2) if they are
diverse enough. We leverage \textit{conditional} models to generate samples of known digit
identity. Thus, we train an SO(3)-equivariant classifier - with the same architecture as H-(V)AE's
encoder - to classify the digit identity of real data at 97\% accuracy, as well as
\textit{conditional} H-VAE models with varying $\beta$. We then sample the conditional models to
construct synthetic datasets of 10,000 spherical images, balanced across digit identity. Finally,
we compute 1) the accuracy of the classifier on the synthetic datasets, and 2) the empirical
variance of generated samples. We find that there is a trade-off between the two metrics: if
$\beta$ is too low, the generate samples are very diverse but do not resemble real digits, whereas
if $\beta$ is too high, samples of the same digit start looking more and more alike
(Figure~\ref{fig:mnist_conditional_samples_VAE_z_16_beta_2pt0}). We empirically select $\beta =
0.6$ and $\beta = 2.8$ for the $z = 16$ and $z = 120$ models, respectively, based on the trade-off
between the above two metrics.
We note that this procedure relies on the assumption that the same regularization strength $\beta$
has the same effect on conditional vs. un-conditional models on sample quality. We believe this to
be likely for the following three reasons: 1) the reconstruction loss is similar for same values of
$\beta$ (Figure~\ref{fig:cosine_loss_mnist_cond_vs_uncond_ablation_in_beta}); 2) the optimal values
of $\beta$ found using the conditional models approximately optimize the latent space metrics of
un-conditional models (Figure~\ref{fig:ablation_in_beta_uncond}); 3) visual inspection of the
conditional vs. unconditional samples across same values of $\beta$
(Figures~\ref{fig:mnist_conditional_samples_VAE_z_16_beta_0pt6}
and~\ref{fig:mnist_samples_VAE_z_16_beta_0pt6},~\ref{fig:mnist_conditional_samples_VAE_z_120_beta_2pt8} and~\ref{fig:mnist_samples_VAE_z_120_beta_2pt8},~\ref{fig:mnist_conditional_samples_VAE_z_16_beta_2pt0} and~\ref{fig:mnist_samples_VAE_z_16_beta_2pt0}).
We perform this procedure only on the NR/R dataset and later use the same selected $\beta$ to train
models on the R/R dataset.


\subsection{Shrec17}
\label{sec:shrec17_details}
\textbf{Model architectures.} Both AE and VAE models have $z = 40$, 7 blocks, $\text{DegreesList} =
[14, 14, 14, 8, 4, 2, 1]$, $\text{MultList} = [12, 12, 12, 20, 24, 32, 40]$, $C_\text{init} = 12$,
with a total of 518k parameters.\\
\\
\textbf{Training details.} We keep the learning schedule as similar as possible for all models. We
use $\alpha = 1000$. We train all models for $120$ epochs using the Adam optimizer with default
parameters, a batch size of $100$, and an initial learning rate of $0.0025$, which we decrease
exponentially by two orders of magnitude over the entire $120$ epochs. For VAE models, we use
$\beta = 0.2$, $E_\text{rec} = 25$ and $E_\text{warmup} = 10$. We utilize the model with the lowest
loss on validation data, only after the end of the warmup epochs for VAE models. Training took
$\sim 11$ hours on a single NVIDIA A40 GPU for each model.


\subsection{Embedding of amino acid conformations}
\label{sec:toy_aminoacids_details}

\noindent \textbf{Pre-processing of protein structures.} We sample residues from the set of
training structures pre-processed as described in Sec.~\ref{sec:protein_neighborhoods_details}.\\

\noindent \textbf{Fourier projection.} We set the maximum radial frequency to $N = 20$ as  it
corresponds to a radial resolution matching the minimum inter-atomic distances after rescaling the
atomic neighborhoods of radius  $10.0\AA$ to fit within a sphere of radius 1.0, necessary for
Zernike transform.\\
The multiplicity composition of the data tensors can be described  in a notation - analogous to
that used by e3nn but without parity specifications - which specifies the multiplicity $C_{\ell}$
for each feature of degree $\ell$ in single units $C_{\ell}\text{x}\ell$: 44x0 + 40x1 + 40x2 + 36x3
+ 36x4.\\
\\
\textbf{Model architectures.} All models have $z = 2$, 6 blocks, $\text{DegreesList} = [4, 4, 4, 4,
2, 1]$, $\text{MultList} = [60, 40, 24, 16, 16, 8]$, $C_\text{init} = 48$, with a total of 495k
parameters. We note that the initial projection is necessary since the multiplicity differs across
feature degrees in the data tensors.\\
\\
\textbf{Training details.} We keep the learning schedule as similar as possible for all models. We
use $\alpha = 400$. We train all models for $80$ epochs using the Adam optimizer with default
parameters and an initial learning rate of $0.005$, which we decrease exponentially by by one order
of magnitude over $25$ epochs. For VAE models, we use $E_\text{rec} = 25$ and $E_\text{warmup} =
10$. We utilize the model with the lowest loss on validation data, only after the end of the warmup
epochs for VAE models.\\
We vary the batch size according to the size of the training and the validation datasets. We use
the following (dataset\_size-batch\_size) pairs: (400-4), (1,000-10), (2,000-20), (5,000-50),
(20,000-20). Training took $\sim 45$ minutes on a single NVIDIA A40 GPU for each model.\\
\\
\textbf{Evaluation.} We perform our data ablations by considering training and validation datasets
of the following sizes: 400, 1,000, 2,000, 5,000 and 20,000. We keep relative proportions of
residue types even in all datasets. We perform the data ablation with H-AE as well as H-VAE models
with $\beta = 0.025$ and $0.1$.\\
We further perform a $\beta$ ablation using the full (20,000) dataset, over the following choices
of $\beta$: $[0 \text{(AE)}, 0.025, 0.05, 0.1, 0.25, 0.5]$.\\
\\
For robust results, we train 3 versions of each model and compute averages of quantitative metrics
of reconstruction loss and classification accuracy.\\
\\
For a fair comparison across models with varying amounts of training and validation data, we
perform a 5-fold cross-validation-like procedure over the 10k test residues, where the classifier
is trained over 4 folds of the test data and evaluated on the fifth one. If validation data is
needed for model selection (e.g. for LC), we use 10\% of the training data.

\subsection{Embedding of protein structure neighborhoods}
\label{sec:protein_neighborhoods_details}
\noindent \textbf{Pre-processing of protein structures}.
We model protein neighborhoods extracted from tertiary protein structures from the Protein Data
Bank (PDB)~\cite{berman_protein_2000}. 
We use ProteinNet’s splittings for training and validation sets to avoid redundancy, e.g. due to
similarities in homologous protein domains~\cite{alquraishi_proteinnet_2019}. Since PDB ids were
only provided for the training and validation sets, we used ProteinNet’s training set as both our
training and validation set and ProteinNet’s validation set as our testing set. Specifically, we
make a $[80\%, \, 20\%]$ split in the ProteinNet's training data to define our training and
validation sets. This splitting resulted in 10,957 training structures, 2,730 validation
structures, and 212 testing structures. \\ %
\\
\textbf{Projection details.} 
We set the maximum radial frequency to $N = 20$ as it corresponds to a radial resolution matching
the minimum inter-atomic distances after rescaling the atomic neighborhoods of radius  $10.0$ \AA\,
to fit within a sphere of radius 1.0, necessary for Zernike transform. We vary $L$ and construct
models tailored to each one (Table~\ref{table:NBs_model_architectures}). \\
\\
\textbf{Model architectures.} See Table~\ref{table:NBs_model_architectures}. We note that the
initial projection ($C_{\text{init}}$) is necessary since the multiplicity differs across feature
degrees in the data tensors, and the ETP necessitates equal multiplicity for all degrees.\\
\\
\textbf{Training details.} We keep the learning schedule the same for all models. We use $\alpha =
1000$. We train all models for $8$ epochs using the Adam optimizer with default parameters, a batch
size of 512, and a constant learning rate of $0.001$. We utilize the model with the lowest loss on
validation data.\\

\begin{table}[h!]
\begin{center}
\resizebox{1.0\textwidth}{!}{
\begin{tabular}{c c c c c c c c}
    \bf $z$ & \bf $L$ & \bf Tensor size & $\mathbf{C_{init}}$ & \bf MultList & \bf DegreesList &
    \bf \# Params & \bf Training Speed (hours) \\
    \hline
     64 & 3 &  616 & 44 &    [50,50,50,64,64,64] & [3,3,3,3,2,1] & 631k &  3.0 \\
     64 & 4 &  940 & 44 &    [50,50,50,64,64,64] & [4,4,4,4,2,1] & 950k &  3.8 \\
     64 & 6 & 1708 & 44 &    [50,50,50,64,64,64] & [6,6,6,4,2,1] & 1.6M &  5.9 \\
     64 & 8 & 2604 & 44 &    [50,50,50,64,64,64] & [8,8,8,4,2,1] & 2.4M &  9.3 \\
    \hline
    128 & 3 &  616 & 44 &  [60,60,60,90,128,128] & [3,3,3,3,2,1] & 1.2M &  3.8 \\
    128 & 4 &  940 & 44 &  [60,60,60,90,128,128] & [4,4,4,4,2,1] & 1.7M &  4.4 \\
    128 & 6 & 1708 & 44 &  [60,60,60,90,128,128] & [6,6,6,4,2,1] & 2.6M &  6.0 \\
    128 & 8 & 2604 & 44 &  [60,60,60,90,128,128] & [8,8,8,4,2,1] & 3.8M & 10.1 \\
    \hline
    256 & 3 &  616 & 44 & [90,90,90,120,200,256] & [3,3,3,3,2,1] & 2.8M &  4.9 \\
    256 & 4 &  940 & 44 & [90,90,90,120,200,256] & [4,4,4,4,2,1] & 3.8M &  5.6 \\
    256 & 6 & 1708 & 44 & [90,90,90,120,200,256] & [6,6,6,4,2,1] & 4.4M &  6.0 \\
    256 & 8 & 2604 & 44 & [90,90,90,120,200,256] & [8,8,8,4,2,1] & 7.9M & 14.5
\end{tabular}
}
\end{center}
\caption{{\bf Model architecture hyperparameters and training time for each H-AE trained on protein
neighborhoods.}}
\label{table:NBs_model_architectures}
\end{table}




\subsection{Latent space classification}
\label{sec:latent_space_classification}
\noindent \textbf{Linear classifier.}
We implement the linear classifier as a one-layer fully connected neural network with input size
equal to the invariant embedding of size $z$, and output size equal to the number of desired
classes. We use cross entropy loss with logits as training objective, which we minimize for 250
epochs using the Adam optimizer with batch size 100, and initial learning rate of 0.01. We reduce
the learning rate by one order of magnitude every time the loss on validation data stops improving
for 10 epochs (if validation data is not provided, the training data is used). At evaluation time,
we select the class with the highest probability value. We use PyTorch for our implementation.

\noindent\textbf{KNN Classifier}
We use the \texttt{sklearn}~\cite{pedregosa_scikit-learn_2011} implementation with default
parameters. At evaluation time, we select the class with the highest probability value.


\subsection{Clustering metrics}
\label{sec:clustering_metrics}
\noindent\textbf{Purity~\cite{aldenderfer_cluster_1984}.} We first assign a class to each cluster
based on the most prevalent class in it. Purity is computed as the sum of correctly classified
items divided by the total number of items. Purity measures classification accuracy, and ranges
between $0$ (worst) and $1$ (best).\\
\\
\textbf{V-measure~\cite{rosenberg_v-measure_2007}.} This common clustering metric strikes a balance
between favoring homogeneous (high homogeneity score) and complete (high completeness score)
clusters. Clusters are defined as homogeneous when all elements in the same cluster belong to the
same class (akin to a precision). Clusters are   defined as  complete when all elements belonging
to the same class are put in the same cluster (akin to a recall). The V-measure is  computed as the
harmonic mean of homogeneity and completeness in a given clustering.

\subsection{On the complementary nature of classification accuracy and clustering metrics}
The clustering metrics ``purity" and ``V-measure” and the supervised metric  ``classification
accuracy” characterize different qualities of the latent space, and, while partly correlated, they
are complementary to each other.\\
\\
Both classes of metrics are computed by comparing the ground truth labels to the predicted labels,
and they mainly differ by how the predicted labels are assigned; the clustering metrics use  an
unsupervised clustering algorithm, while the classification metric uses a supervised classification
algorithm to do so. As a result, these metrics focus on different features of the latent space. For
example, the clustering metrics are largest when the test data naturally forms clusters with all
data points of the same label.  While this case can result in a high supervised classification
accuracy, clustering is not a necessary condition for high classification accuracy. Indeed, the
supervised signal could make the predicted labels depend more heavily on a subset of the latent
space features, instead of relying on all of them equally, which is what the clustering algorithm
naturally does. Therefore, it is reasonable to conclude that having higher clustering metrics and a
lower classification accuracy is a sign that class-related information is more evenly distributed
across the latent space dimensions. Overall, the complementary aspect of these metrics makes it
necessary to use all of them when comparing the performance of different models in each task.


\section{Protein-Ligand Binding Affinity Prediction}
\label{sec:lba_details}
\subsection{Data Preprocessing Details}
We leverage the H-AE models trained on Protein Neighborhoods
(Sec.~\ref{sec:protein_neighborhoods_details}) to predict the binding affinity between a protein
and a ligand, given their  structure complex-- an important task in structural biology. For this
task, we use the data from the ``refined-set" of PDBBind~\cite{su_comparative_2019}, containing
$\sim 5,000$ structures. We use the dataset splits provided by ATOM3D~\cite{townshend_atom3d_2021}
to benchmark our predictions. ATOM3D provides two splits based on the maximum sequence similarity
between proteins in the training set and the validation / test sets. We use the most challenging
split for our benchmarking in which the sequences of the two sets have at most 30\% similarity.

For each complex, we first identify residues in the binding pocket, which we define as residues for
which the C$\alpha$ is within 10\AA\, of any of the ligand's atoms. We then extract the 10\AA\,
neighborhood for each of the pocket residues that can contain atoms from both the protein and the
ligand. We compute the ZFT (Eq.~\ref{eq:ZFT}) for each neighborhood with maximum degree $L$
matching the degree used by the H-AE of interest. We then compute residue-level SO(3)-invariant
embeddings by running the pre-trained H-AE encoder on the neighborhoods surrounding each of the
residues in the pocket. We sum over the residue-level embeddings to compute a single pocket-level
embedding, which is SE(3)-invariant by construction. Each pocket embedding is used as a feature
vector within a standard Machine Learning Regression model to predict protein-ligand binding
affinity, provided by PDBBind either in the form of a \textit{dissociation constant} $K_d$, or an
\textit{inhibition constant} $K_i$. Due to data limitations, and consistent with other studies on
LBA, we do not distinguish between $K_d$ and $K_i$ and regress over  the negative log of either of
the constants that is provided by PDBBind; the transformed quantity is closely related to the
binding free energy of protein-ligand interactions. 

The protein-ligand binding free energy is an extensive quantity, meaning that its magnitude depends
on the number of residues in the binding pocket. In other words, a larger protein-ligand complex
can establish a stronger binding. The pocket embedding, which we define as the sum of the
residue-level embeddings, is a simple quantity that is extensive and SE(3)-invariant, and
therefore, suitable for regressing over protein-ligand binding free energy. Nonetheless, this map
is not unique and other extensive transformations to pool together residue-level embeddings into a
pocket-level embedding can be used for this purpose.

It should be noted that in the protein-ligand structural neighborhoods we only include the ligand
atoms that are found in proteins and were used in the training of the H-AE models (i.e., C, N, O,
and S). About 31\% of the ligands contain other kinds of atoms, and since our model does not ``see"
these atoms, we hypothesize that our predictions are worse in these cases. In fact, for H-AE+R.F.
(Table~\ref{table:lba_results}), the Spearman's correlation over the ligands containing only
protein atoms is higher than for the ligands containing other kinds of atoms by about 0.03 points.
Therefore, we expect training H-(V)AE to recognize more atom types - or at  least making it aware
that some other unspecified non-protein atom is present - would yield even better results; as
protein structures are often found in complex with other non-protein entities, e.g. ligands and
ions, there is training data available for constructing such models.


\subsection{Machine Learning models used for prediction}
\noindent \textbf{Linear Regression.} We implement a linear regression model as a one-layer fully
connected neural network with one output layer. We use MSE loss as training objective, which we
minimize for 250 epochs using the Adam optimizer with batch size 32, and initial learning rate of
0.01. We reduce the learning rate by one order of magnitude every time the loss on validation data
stops improving for 10 epochs. We use PyTorch for our implementation.\\

\noindent \textbf{Random Forest Regression.} We use the
\texttt{sklearn}~\cite{pedregosa_scikit-learn_2011} implementation. We tune hyperparameters via
grid search, choosing the combination minimizing RMSD on validation data. We consider the following
grid of hyperparameter values:

\begin{itemize}
    \item max\_features: [$1.0$, $0.333$, sqrt]
    \item min\_samples\_leaf: [$2$, $5$, $10$, $15$]
    \item n\_estimators: [$32$, $64$, $100$]
\end{itemize}

\subsection{Extended Discussion of Results}

Table~\ref{table:lba_results_appendix} shows benchmarking results on the split of ATOM3D with 30\%
sequence similarity, extending Table~\ref{table:lba_results} in the main text with more baselines.
Notably, we show the prediction results for the models that only use the SO(3)-invariant ($\ell =
0$) component of each neighborhood's Zernike transform (our Zernike inv. baseline), instead of the
complete embeddings learned by H-AE. Using H-AE consistently outperforms this baseline, implying
that invariant information encoded in higher spherical degrees, and extracted by H-AE, can lead to
more expressive models for downstream regression tasks.

